\documentclass[12pt]{letter}
\usepackage[margin=1in]{geometry}
\usepackage{hyperref}

\signature{Humza Hafeez\\Independent Researcher\\ORCID: 0009-0000-0853-8807}
\address{humza@closure-theory.org}

\begin{document}

\begin{letter}{The Editors\\Physical Review Letters}

\opening{Dear Editors,}

We submit for your consideration the manuscript ``Geometric Non-Stationarity of the Type~Ia Supernova Standardization Manifold and Its Impact on Dark Energy Inference.''

This work demonstrates that the covariance structure of the SN~Ia standardization space $(x_1, c, \Delta\mu)$ is strongly redshift-dependent, using three independent statistical tests on the Pantheon+ compilation of 1,590 supernovae:

\begin{enumerate}
\item \textbf{Box's M test} ($p = 2.3 \times 10^{-7}$): the covariance matrices of the standardization space differ across redshift groups at high significance.
\item \textbf{Rolling channel--residual correlations} ($p = 3 \times 10^{-30}$): the coupling between color and Hubble residuals reverses sign from $+0.09$ at low $z$ to $-0.44$ at high $z$.
\item \textbf{Interaction terms} ($\Delta\chi^2 = 269$ for 2 parameters): adding $z \times x_1$ and $z \times c$ to the Tripp equation overwhelmingly rejects fixed-basis standardization.
\end{enumerate}

These results establish that the standard Tripp equation---the foundation of all current SN-based dark energy constraints---applies a fixed projection to a nonstationary manifold, producing structured residuals that mimic cosmological signals.

The immediate consequence is significant: when the $w_0 w_a$ dark energy parameterization is refit under corrected standardization, the best-fit $w_a$ shifts from $-1.88$ to $-0.02$. $\Lambda$CDM with nonstationary Tripp outperforms $w_0 w_a$ with standard Tripp by $\Delta\chi^2 = 37$. This directly impacts the interpretation of recent DESI BAO results reporting evidence for time-varying dark energy ($w_a = -0.75 \pm 0.29$; arXiv:2404.03002).

We believe this result is appropriate for Physical Review Letters because:

\begin{itemize}
\item It identifies a previously unrecognized geometric systematic in SN~Ia standardization that affects every major cosmological analysis using Pantheon+, DES-SN5YR, or DESI SN anchors.
\item The proof is mathematical (linear algebra and standard multivariate statistics), not model-dependent.
\item Analogous channel-selective degradation in SDSS quasars and CHIME fast radio bursts suggests the effect is not unique to supernovae.
\item The result has immediate practical implications for ongoing dark energy programs (DESI, Euclid, Rubin/LSST).
\end{itemize}

All analysis scripts and data products are publicly available at \url{https://github.com/Chaiwalah/closure-theory}. The manuscript has not been submitted elsewhere.

We suggest the following referees with relevant expertise in SN~Ia standardization and cosmological inference:

\begin{itemize}
\item Dan Scolnic (Duke University) --- Pantheon+ lead
\item Dillon Brout (Boston University) --- Pantheon+ cosmology lead
\item Tamara Davis (University of Queensland) --- SN~Ia standardization systematics
\item Ryan Keeley (KASI) --- Pantheon+ covariance analysis
\end{itemize}

\closing{Sincerely,}

\end{letter}
\end{document}
